\documentclass{ximera}


%\usepackage{currfile}

%\makeatletter
%\ifxake
% The code below the \else is executed in a sagecell on the Ximera
% server, so \makerandom doesn't have to do anything when run under
% xake.
%\newcommand{\makerandom}{}
%\else
%\newcommand{\makerandom}{%
%  \ST@wsf{jobname="\currfilebase"}%
%  \ST@wsf{import hashlib}%
%  %\ST@wsf{set_random_seed(int(hashlib.sha256(jobname.encode('utf-8')).hexdiges%t(), 16))}%
%}
%\fi
%\makeatother


\title[Breal-grond]{This is the sample! Waffle!  Iff!}

\begin{document}
\begin{abstract}
  The begining and now an instructor! I am doing math like $\arcsin$. AGAIN
\end{abstract}

\maketitle
%\makerandom

HELLO.

compare $\mathrm{x}$ to x.

compare $M$ to \textit{M}.

But also if $y^2 = 4y$ then if $y \neq 0$ we may compute $\cancel{y} \cdot y = 4 \cancel{y}$ so $y = 4$.

$21 \cancel{\textrm{ feet}}$

Look at this:
\[
  \pm 2
\]

%I roll a D20 and I get $\sage{randint(1,20)}$.

%I roll a D1000 and I get $\sage{randint(1,1000)}$.

\marginpar{Pretty fun huh?}

%\begin{sagesilent}
%  lhs = randint(0,10)
%  rhs = randint(0,10)
%  x = rhs - lhs
%\end{sagesilent}

%I can concatenate with $cat + dog = \sagestr{'cat' + 'dog'}$.

%\begin{problem}
%  Suppose $\sage{lhs} + x = \sage{rhs}$.

%  Then $x = \answer{\sage{x}}$.
%\end{problem}

It is a simple exercise to solve the system of two equations
\begin{equation} \label{small}
\begin{array}{rcrcr}
 x & + & y & = & 7 \\
-x & + & 3y & = & 1
\end{array}
\end{equation}
to find that $x=5$ and $y=2$.  One way to solve
system \eqref{small} is to add the two equations, obtaining
\[
4y=8;
\]
hence $y=2$.  Substituting $y=2$ into the $1^{st}$ equation in
\ref{small} yields $x=5$.

\begin{enumerate}
\item\label{item:one} The first
\item The second
\item\label{item:three} The third
\end{enumerate}

\begin{problem}
  Also testing \wordChoice{\choice{incorrect}\choice[correct]{yas}\choice{$2x$}\choice{so wrong}} and such.
\end{problem}


I am testing the waffle and the iff statement.  But consider \ref{item:three} or that first \ref{item:one}.

\begin{image}
  \begin{tikzpicture}
    %\draw (0,0) circle (1in);
    %\draw (0,0) circle (0.25in);
    \draw (0,0) -- (1in,1in);
    \draw (1in,0) -- (0,1in);
  \end{tikzpicture}
  \alt{these are circles}
\end{image}

\begin{image}
  \begin{tikzpicture}
    \draw (0,0) -- (1in,1in);
  \end{tikzpicture}
  \alt{this is a line segment}
\end{image}


\begin{theorem}
  \label{thm:blah}
\end{theorem}

We might look at \ref{thm:another} but consider \ref{thm:whee} or even \ref{thm:blah}.

\begin{theorem}
  \label{thm:whee}  Whee theorem.
\end{theorem}
\begin{proof}
This is the proof for reals.
\end{proof}

\end{document}


% Changed at Tue Jul 28 00:06:03 EDT 2015
% Changed at Tue Jul 28 00:06:04 EDT 2015
% Changed at Tue Jul 28 00:06:28 EDT 2015
